\section{A Non-Linear Interlude: Counting Register Shifts in a Scene Arc}\label{sec:markov-interlude}

Before the main theorems, a short example worth walking through.
Section~\ref{sec:mergeable-sketches} gave the classical recipe for a tree
of summaries: encode shards, merge bounded states, and read out an
aggregate query. Section~\ref{sec:framework} generalized that recipe to a
learned summarizer under three local laws. The simplest task that shows
why the learned generalization matters is one whose oracle is an aggregate
but cannot be computed by an unordered register-wise merge unless the
state carries boundary information. We walk through that task here as a
scene-register arc illustrated through \emph{Hamlet}, and motivate why
$g$ must be a richer operator than a set-like sketch merge.

\subsection{The Running Example: A Scene-Register Arc}\label{sec:scene-register-example}

Consider one judgment about a play's emotional shape: how many times
does the dominant emotional register change from one scene to the
next? Computing this by hand requires reading every scene, deciding
which register it sits in, and walking the sequence of decisions
end-to-end. We take \emph{Hamlet}, partition it at the scene level,
and assign each scene a single dominant register from the
seven-element alphabet
\[
  \RegSet \;=\; \{\text{Mourning},\, \text{Foreboding},\, \text{Political},\, \text{Intrigue},\, \text{Madness},\, \text{Action},\, \text{Comic}\}.
\]
The boundary between 3.3 and 3.4 (Claudius at prayer to Polonius killed)
flips from Intrigue to Action, and the boundary between 4.7 and 5.1 (the
poisoned-foil plot to the gravediggers) flips from Intrigue to Comic.
These two flips give a concrete picture of what the oracle counts.

The oracle $\fstar$ counts register shifts across adjacent scene
boundaries. It is an integer aggregate of the same shape as a classical
sketch query, but its value depends on an \emph{ordered} pair of
neighbouring states, not on an unordered bag of features. That is why
the example matters: the moment we split the play into two halves and
summarize each half independently, we risk losing the pair of registers
on either side of the split point.

\subsection{What We Assume and What Is Illustrative}\label{sec:scene-register-assumptions}

Three things are load-bearing in this example, and they are worth
distinguishing.

First, we assume that each scene has a \emph{dominant register} (a
single label that best describes its emotional tone), and that the
register can shift across adjacent scenes. This is the structural
assumption: scenes are labelable, labels live in a finite alphabet, and
the oracle is count-of-shifts. Everything the framework proves about
this example holds as long as that structural assumption holds.

Second, we assume a specific alphabet $\RegSet$ and a specific labeling
of Hamlet scenes into that alphabet. This is \emph{illustrative}. A
reader who disagrees with ``Mourning'' or thinks Act~1 Scene~2 is
Foreboding rather than Political is doing literary criticism; the
framework is indifferent to the answer, because the mechanism (leaf
summaries must carry enough information to reconstruct the boundary
registers) does not depend on which specific register any scene gets.
We chose seven registers because seven is large enough to be a real
categorical distinction and small enough to be readable. Four would
work. Twelve would work. The only thing that would not work is an
alphabet so fine-grained that adjacent scenes almost always differ, or
so coarse that they almost never do.

Third, we do not claim to have coded Hamlet. The empirical results in
Section~\ref{sec:markov-walkthrough} come from a \emph{synthetic
data-generating process} that produces sequences with the same shift
structure as this example but uses neutral register labels.
\emph{Hamlet} is a narrative container; the DGP is what the MAE
curves actually measure. The example shows a task whose oracle defeats
unordered register-wise merges, and that pedagogical purpose does not
require a literary-critical commitment.

\subsection{The Synthetic Data-Generating Process}\label{sec:synthetic-dgp}

The DGP is a Markov chain on $K$ hidden registers whose transition
matrix is chosen so that adjacent states differ about as often as
adjacent scenes differ in the Hamlet illustration. A generated document
is a sequence of $128$ positions, each emitting one observed token from
a small palette per register. The palettes are disjoint (no token
belongs to two registers), but the learner does not know which tokens
share a register. It must discover the groupings from data. The oracle $\fstar$ counts the positions at which the hidden
register changes. A typical document has about five shifts;
Appendix~\ref{app:mechanism-checks} gives the full specification,
including the $K$-register variants used for the hard-case regime and
the per-seed sample sizes underlying every reported curve.

We refer to positions here as ``tokens'' and to the hidden state as a
``register'' because that terminology matches the Fourier-neural-operator
architecture we will plug in at Section~\ref{sec:fno-primer}. The
Hamlet framing of \S\ref{sec:scene-register-example} is the same
mechanism viewed at scene granularity: positions become scenes, the
hidden state becomes a dominant emotional register, and the oracle
becomes the count of register shifts. The structural claim that an
unordered leaf summary erases boundary information does not care
whether the unit is a token or a scene.

\begin{figure}[t]
    \centering
    \includegraphics[width=\linewidth]{markov_scene_arc_hamlet.pdf}
    \caption{\textbf{An illustrative scene-register arc for \emph{Hamlet}.} Consecutive scenes are arrayed left to right; colors encode the dominant register assigned to each scene under the seven-register alphabet $\RegSet$. The canonical boundary flips at Acts 3.3$\to$3.4 and 4.7$\to$5.1 motivate why leaf summaries must carry boundary state. The register assignments are an illustrative mood taxonomy, not a literary-critical claim; the empirical results come from the synthetic data-generating process of~\S\ref{sec:synthetic-dgp}, not from annotated Hamlet text.}
    \label{fig:hamlet-registers}
\end{figure}

\subsection{What Goes Wrong with an Unordered Merge}

Suppose we attempt to apply the HyperLogLog recipe directly: encode each
leaf into a register-like summary, merge pairwise with a commutative
operator, read out the count. For this to succeed the encode step would
have to reduce each leaf to a feature set that, under commutative merge,
still distinguishes between (i) a document whose two halves happen to end
and begin with the same register and (ii) a document whose two halves end
and begin with different registers. No register-wise max, additive count,
or bitmap-OR merge can distinguish these two cases: both produce the
same leaf-level aggregate. The only distinguishing information is
\emph{ordered and positional}: which register is at the right edge of
the left leaf, and which register is at the left edge of the right
leaf. A commutative merge by definition cannot preserve such
information.

The hand-designed sketch that makes the example work carries two pieces
of boundary metadata in addition to the internal shift count: the first
register and the last register of the span the leaf covers. The merge
is then
\[
\text{merge}(L, R) \;=\; \Big(L.\text{count} + R.\text{count} + \One\{L.\text{last} \neq R.\text{first}\},\; L.\text{first},\; R.\text{last}\Big),
\]
which is associative but \emph{not} commutative. This works, but it
required a human to recognize that the oracle depends on ordered
boundaries and to hand-engineer a state that preserves them. For the
political-science application of Section~\ref{sec:manifesto-llm}, the
analogous engineering for RILE scoring would have to carry forward
every span-level cue the rubric touches (a specification no one has
written down). We instead let $g$ \emph{learn} what boundary
information to carry.

This hand-designed state is also part of the Lean artifact. The
formalization defines the type \texttt{MarkovCountSketch} with cases
\texttt{empty} and \texttt{nonempty(count, first, last)}, proves the
merge operation is a monoid, proves that the exact point-mass
summarizer satisfies the leaf and merge local laws, and invokes the
general preservation theorem to obtain zero root distortion. The same
file shows that C2 is independent of leaf and merge correctness: a
simple on-range ``flip'' summarizer satisfies the one-pass laws but
breaks under re-summarization. Appendix~\ref{app:proof-artifacts}
records the corresponding theorem names.

We instantiate $g$ as a Fourier neural operator (a position-aware architecture whose frequency-domain mixing carries boundary context across the full sequence window). That architecture supplies representational capacity, not certification by itself. The certification target is still the local-law subspace: C1 asks each leaf to keep the internal count and boundary registers it needs, and C3 asks the merge to preserve the ordered boundary correction. Section~\ref{sec:fno-primer} introduces the realizability layer. Section~\ref{sec:markov-walkthrough} then shows that a tree of FNO leaves solves this task; Section~\ref{sec:hll-parity} shows the same framework recovering the classical HyperLogLog sketch when the oracle is a register-friendly query.
